\section{Spiritual Challenge and Renewal}

\begin{quotationx}
A religious factor is necessary as the background for a true heroic conception of life that must be essential for our
political alignment. It is necessary to feel in oneself the evidence that there is a higher life beyond this
terrestrial life, because only those who feel this way possess an unbreakable and unconquerable strength, only they
will be capable of an absolute enthusiasm. \flright{\textsc{Julius Evola}} 

\end{quotationx}
After the devastation of World War II, \textbf{Julius Evola} was asked to formulate a program for the political and
spiritual renewal of post war Europe. Evola was able to recognize the ideal of highest spiritual tradition of Europe,
which he described this way:

\begin{quotationx}
Certainly, if Catholicism were capable of making a program of high ascesis its own and exactly on this base, almost like
a recovery of the spirit of the best Medieval crusader, makes of faith the soul of an armed bloc of forces, almost a
new close-knit Templar Order, relentless against the currents of chaos, breakdowns, subversion, and the practical
materialism of the modern world — certainly, in such a case, and also in the case that, at the minimum, it
held firm to the positions of the Syllabus, there could not be a single instance of doubt about our choice. \flright{\textsc{Julius
Evola}, \textit{Point 11}\footnote{\url{https://www.gornahoor.net/?p=4828}}} 

\end{quotationx}
Unfortunately, after looking around — rather cursorily in our opinion — he determined that there
was currently no organisation that supported such a spiritual orientation, as it perhaps had done in the past. Hence,
he concluded that a “pure reference to a transcendental spirit” would suffice. Well, more than six decades later, it
certainly has not sufficed.

We can start, for example, with the Syllabus\footnote{\url{https://www.papalencyclicals.net/Pius09/p9syll.htm}} mentioned, which is the 19th century Syllabus of Errors. It is an easy to
read collection of false opinions. Who today would deny all, or at least most, of those errors? We don't
see that even among \emph{soi disant} Traditionalists. Evola's attitude was that if George
won't do it, then he'll go his own way. And he took many along with him. Despite his vast
erudition, Evola always remained at the level of doxa (opinions) and dianoia (theories). He never — his
writings and interviews show this — was able to reach the level of episteme, the goal of all valid
Traditions. Obviously, a vague affirmation of some transcendental spirit goes nowhere.

\paragraph{The True Revival}
In point of fact, there has been, and still is, such a program. Although not centrally organised, there are texts and
even small groups working in the USA and in Europe, that fulfil the conditions mentioned above. That is why, for us,
there is no doubt about our choice.

For example, \textbf{Wolfgang Smith}, starting from the insights of Rene Guenon, has revived such a tradition within
Catholicism. Specifically, he has written on the Christian Kabbalah, the chakras, different levels of existence, Jacob
Boehme, Meister Eckhart, in an effort to bring more depth into the Western tradition.

\textbf{Boris Mouravieff}, working completely independently, has revealed and reformulated the teachings of a monastery
on Mount Athos. A glance at his bibliography shows many familiar names, including Boehme, Eckhart, and Guenon.

But \textbf{Valentin Tomberg} has gone the furthest, particularly in his \emph{Meditations on the Tarot}, which is the
summa of his life's work. First of all, he addresses Evola's concern incidentally. As a
spiritual reality, the Catholic Church is the Bride of Christ. However, as a human institution, it also has an
egregore, a phantom being, which is a human-created parody of the spiritual reality. Many people get fixated on the
egregore and reject the whole tradition. Those with spiritual vision, on the other hand, see a living stream of truth
and life, as well as a way to achieve them.

\paragraph{An Incomplete Path}
The Meditations are a massive accomplishment. Nevertheless, there is much that is left out, and he expects the task that
he initiated to be continued by those he considers his “friends”. I know some of those who have indeed continued the
work of Christian Hermetism, although not in such a public way as we have. Still, there is much more to be done.

Tomberg left us a long list of precursors. These include philosophers, saints, Hermetists, theologians, scientists,
psychologists, inter alia. We have endeavoured to bring many of their writings to a wider attention. However, as a
precursor to some planned articles, we need to make here some specific points. In particular, we need to put his
earlier writings into context with his later. Although he allegedly desired to be known by the Meditations, there are
some allusions that make that impossible.

\paragraph{Relationship to Steiner}
Although Tomberg became quite critical of the Anthroposophical Society (and vice versa), he always held \textbf{Rudolf
Steiner} in high regard. For example, he wrote this in Letter XV:

\begin{quotationx}
And yet Rudolf Steiner has certainly said things of a nature to awaken the greatest creative elan! His series of
lectures on the four Gospels, his lectures at Helsingfors and Dusseldorf on the celestial hierarchies — without mentioning his book on the inner work leading to initiation (\emph{Knowledge of the Higher Worlds and its
Attainment}) — would alone suffice to inflame a deep and mature creative enthusiasm in every soul who
aspires to authentic experience of the spiritual world. 

\end{quotationx}
Nevertheless, there is really nothing in the Meditations specifically referencing those works. I've read
most of those lectures, and there are things in them that may be quite surprising. Nevertheless, we have to assume that
they play a background role in the Meditations, even if not explicitly stated. This will be confirmed in the next
section.

\paragraph{Cosmology}
A revived tradition requires a suitable cosmology, certainly one that can challenge the generally accepted scientific
version. In \textit{Letter~X}, we are given some sources for the creation of such a cosmology.

\begin{quotationx}
Firstly, the Fall... here we are confronted with the Biblical account of paradise and the six days of creation; with
the impressive tableau of natural evolution that science advances; with the contours of a majestic outline by the
genius of ancient India of kalpas, manvantaras and yugas — a world of periodicity and rhythm, a world
dreamt periodically by cosmic consciousness; with the exposition (following the \emph{Stanzas of Dzyan}) of cosmogony
and anthropogony according to the Indo-Tibetan tradition, given by H. P. Blavatsky in the three volumes of her
\emph{Secret Doctrine}; with the grandiose tableau of the spiritual evolution of the world through seven so-called
“planetary” phases that Rudolf Steiner has bequeathed to the dumbfounded intellectuality of our century; lastly, with
the cosmogonies and eschatologies — explicit or implicit — of Hermes Trismegistus, Plato, the
Zohar and diverse gnostic schools of the first centuries of our era. 

\end{quotationx}
That is a vast amount of material to master. Tomberg himself acknowledges this:

\begin{quotationx}
May I be permitted to say straight away that, although I have had actual experience of comparing the whole range of
these ideas and documents for more than forty years, I cannot make use of them here in the sense of the treatment which
they merit, i.e. to classify them, to extract the essential points of similarity or contrast, to make relevant
quotations, etc. If I were to do so, I would drown the essential theme in a sea of secondary elements (secondary with
regard to the main theme). Therefore I have to proceed in the following way: the spirit of all the various ideas and
documents enumerated above will be present as a general background, but it will be necessary to refrain from any
explicit use of the material which they comprise. 

\end{quotationx}
In particular, this general background includes Steiner's \emph{Outline of Occult Science}, which Tomberg
alluded to. And this brings the entire story of the spiritual evolution of man, his various bodies, and the activities
of the celestial hierarchies, both left and right, in the formation. These notions will be used in future articles.

\paragraph{Psychology}
Conventional Thomism teaches that man is composed of several sheaths: a physical body, a vegetable soul, an animal soul,
and an intellectual soul. Moreover, it is the higher sheaths that form the lower, which is contrary to conventional
scientific thought. Unfortunately, philosophers are content to leave it at that, and don't bother to
explore the actual consequences of such a teaching.

The esoteric teaching, on the other hand, does so. It may use different names, viz., the physical, etheric, and astral
bodies, but the essential elements are the same. The Hermetist, however, will make the efforts to observe the bodies or
sheaths. In this way, he can also discern the influences of alien elements on them. That is not the end of the story,
since the chakras are also part of one's interiority. The Hermetist will learn about the proper development
of the chakras.

The main challenge is that the effort of the I to form the lower sheaths is not yet completed. There is resistance to
that; besides, the awareness of the I is mostly dim and the will is weak.

The details cannot all be made public, but it will be shared with our groups at the proper time.

\paragraph{Finis}
With this being said, we can next complete the discussion of the three temptations in the wilderness.

\flright{\itshape Posted on 2018-09-09 by Cologero}

\begin{center}* * *\end{center}

\begin{footnotesize}\begin{sffamily}
\texttt{Han Fei on 2018-09-10 at 15:56 said: }

I wish to ask a question. Would the basic points of adherence required of a Christian be enough for the purpose of
personal reintegration, or does the traditional church require the introduction of some sort of esoteric doctrines,
morphed over centuries and even taken from exterior traditions such as Judaism and Hinduism? By basic points I mean
regular attendance of mass, partaking of the sacraments, and obedience of the commandments. As Evola put it, to follow
the Syllabus, I think he meant more than the 19th century document, but rather the fundamental tenets and duties of the
faith.


\hfill

\texttt{Mikkel on 2018-09-11 at 21:46 said: }

Any points of adherence to anything can never be enough for the personal reintegration you speak of, forgive me if I am
miss-assuming the direction of your question. Reintegration sounds to be something one would call more true than simply
surface level action that we see today.

Many Catholics follow what some refer to as “Traditional Catholicism” and many at least follow the basic tenants of the
faith, as referred to in the article, yet we are where we are today in the current state of things, so in short, no it
is not enough. As to the Syllabus, it is being pointed to because of how stark in contrast the Church and anyone else
compares points to the statements being made in it, but I believe your thought on how it is more than the document that
JE refers to is correct. Maybe it would be better to ask what kind of mindset creates the Syllabus, especially if
organic, and how does that mindset reflect the principles enacted (the source)? 

The basic points must be followed but that is not the real action, the right action, behind the fundamental tenants and
duties and what is incorporated as the body of the Church and the world of the people, more so closer I think to the
mindset which I mentioned. The Church does contain truth regardless of how much people wish to distance themselves from
it and look at it with disdain and no matter how the failures of the members and leaders of said church affect others
it will contain it. No one, not one thing or person or excuse or motivation will excuse the lack of adherence to truth
for anyone and the truth of heroic life however will always be within the Church regardless of what anyone does or says
even within our time.

Though our cycle is referred to as the “Kali Yuga”, does this mean that truth somehow is gone now? That the validity of
any sound tradition is gone? It reflects our manifestation in this world, “Ignorance of dharma will occur” does not
mean that dharma is suddenly gone, but the ignorance has appeared. No matter how “pointless” it all is, no ladies
wearing veils, kneeling for the Eucharist, or others chanting in Latin will somehow make that easier for “George” to be
reintegrated or live in a heroic manner. But I do think that investigating the sources of something like the Syllabus
would lead us beyond “doxa (opinions) and dianoia (theories)”.

As Guenon states, an elite is required for these `esoteric' doctrines to be put out into the
world, especially in our age. These doctrines need not be taken from others, since they are already reflected and
perhaps transposed, even if the Church was lacking in such a doctrine, it should be understood from the source. A
variation on a variation leaves much wanting. Maybe this could be done much like seeking the interior of man to
understand, the Church must be understood from within or from the source, it has not lost any of its esoteric
doctrines, only we need to remember better to then enact it. Tomberg found that all of his thought corresponded with
the what the Church taught, perhaps as the article mentions our way of understanding what the Church is trying to say
is only “incomplete” so then it is up to us to complete it.


\hfill

\texttt{argusandphoenix on 2018-09-15 at 23:53 said: }

We're living through the collapse of a civilization, and not just a culture. So the Church has gotten
tangled up in some of the wreckage, or even contributed to it, at times, by dereliction. So the usual strategies (just
go to Church, or be a good citizen) aren't going to necessarily place us in the context of the role we are
meant to play. Unless you are fortunate, you'll have to go and dig, to some degree, and find deeper
currents within and connected to the Church, and in elective affinity or sympathy with it. There's plenty
of that material that C. Salvo has been unearthing or dusting off or pointing out on the website. If you just want to
be saved and live a good life, the Church certainly still makes that possible. But there are specifically Christian
groups of esotericists out there, if that is more of a comfort zone when dealing with “re-integration”.


\hfill

\texttt{Agnostia on 2020-09-09 at 23:30 said: }

I am not at all convinced that a traditionalist re-invigoration of Catholicism, even within the exoteric domain, is
outside the realm of possibility. What is lacking is vision.

It goes without saying that the Church has had little difficulty in incorporating foreign doctrines which are compatible
with Catholic dogma or pre-existing doctrine. Perhaps my favorite example of such is that doctrine of
`guardian angels', which has it's own specific (let's call it)
`parallel' in the `personal daimons' of Apuleius (see On the God of
Socrates, by the same).

It should also be noted that the process of `perfection’, ‘divinization’, ‘theosis' are not at all unlike the procurement of the `diamond-thunderbolt body’ or the discovery of the `stone of the philosophers', the `VITRIOL'ic
process, etc. when dripped of their otherwise heretical, `magical' dross. One such funny
example of this can be seen in the friar Roger Bacon, who in his “On the Non-Existence of Magic”, trashes any and all
notions of sigils, talismans, spells, theurgical and thaumaturgical ceremonial operations, astrological influences and
the like, before giving a near-perfect summation of the hermetic Opus, only half-allegorized as a recipe for obtaining
gun-powder, all the while upholding himself as an orthodox Catholic and material scientist.

What is so great about hermeticism in particular, is that so many of it's key operations can easily be
interpreted, re-contextualized and integrated within a Catholic framework. In fact to do so is hardly new. The
concepts, `rituals’, and steps mentioned in such works as J. Evola's Hermetic
Tradition \& Intro. to Magic series (properly understood, as these require much review), Guenon's
Perspectives on Initiation, and the alchemical texts themselves(my favorite being M. Maiers' Atalanta
Fugiens, but even in the more `magical' grimoires like those of Agrippa or the Picatrix, where
one finds simultaneously tables of magical correspondences alongside dissertations on the `virtue of
religion’!) have their almost direct-correlatives in established and orthodox Christian tradition. With
some research, and by dropping the superficial hermetic conceptual labeling, it's not hard to find
analogous practices in the Philokalia and the other writings of the mystics and Doctors of the Church. Among the
certain circles today, there is also a re-awakening of a sense of an objective `esoteric’ meaning in the Biblical texts which are surprisingly neither anti-hermetic nor outside the purview of established
Orthodoxy.

What must be done is a selection and compilation of these various techniques, yet presented within a new
`Active' framework, for a particular type of individual.

What about the sorts of mysticism and metaphysics common to the Medieval and Renaissance mystics? Can anything be done
with these? There is certainly room for a Platonic-Augustinian Idealism. Recent Popes have affirmed this, among other
ways, in their canonization of St. Hildegard and their appraisal of Pseudo-Dionysius. This is to say nothing of some of
the honourable, if occasionally anti-traditional and often misunderstood intentions, of the 20th century pioneers of
Vatican II.

What about the limitations of the `mystical-devotional attitude' to which Evola so frequently
lamented? This is, in my opinion, the hardest thing to grapple with. While there is obviously nothing inherently wrong
with a purely devotional mindset, traditional Christianity, in my opinion, most certainly has been limited by this
approach, which as currently practised transparently lacks the autonomous impetus desirable for the
`differentiated man’. It is this approach which has also found itself constricted within the
paradigms of a strict Thomism and a rigid pre-occupation with ethics. Undoubtedly this is where the rubber hits the
road, and will take the most effort. But to successfully tackle this phenomenon could have such an affect, as to
re-vivify that spirit which had previously been crushed by that most famous declaration of Nietszche, while
simultaneously affirming the ideal which Nietszche himself placed forth: that of the romantic `free
spirit’.

To aggregate and re-present these things, under an old-yet-new framework, has the potential to change the world. More
importantly, it has the potential to do so within the range of an established institution and a traditional order.
Perhaps there can be no greater `holding action' than the implementation of this, no greater
arsenal of which to use against the `quantitative' and `malefic' forces
of this end of an era.


\hfill

\texttt{Santiago on 2022-03-10 at 21:36 said: }

“[Evola] never... was able to reach the level of episteme.”

I've been thinking about this line.

Compared to Guenon, Evola (“despite his vast erudition”) does seem far less interested in metaphysics, and far more
interested in methods, such as Tantric and yogic practices, cultivation of the chakras, empowerment of the will,
visualization practices, and all manner of other techniques. Guenon, on the other hand, will
discuss these, but always with the emphasis being on their place within a larger framework and a supreme goal. 

It almost seems as if Evola was more interested in prolonging his individual existence, rather than in transcending it
entirely. Is this a fair criticism? And whether or not it is, to what extent do the two paths converge?

\end{sffamily}\end{footnotesize}
